\section{Introduction}

Goal-oriented dialogue agents are increasingly expected to act on behalf of
users, not simply to generate fluent responses. In buyer-side negotiation, for
example, a useful agent must pursue a favorable price while keeping the seller
engaged, reaching an agreement within a limited number of turns, and respecting
the buyer's budget. Target-driven recommendation involves a similar balancing
act: the agent should guide the user toward a desired item without degrading
satisfaction or conversational efficiency. These objectives are not independent.
An aggressive bargaining move may improve price utility if the seller remains
cooperative, yet the same move can reduce the probability of closing the deal.
The policy therefore has to reason about strategic trade-offs and constraints
as part of the dialogue process itself.

The difficulty is compounded by the way goals are specified in practice.
Traditional goal-oriented policies are usually trained for a fixed objective,
which makes adaptation to new user requirements expensive. Recent
multi-objective methods reduce this cost by conditioning a policy on preference
vectors, but the interface remains rigid: it assumes that the user's goal has
already been translated into explicit numerical weights. In realistic
deployment, users are more likely to provide natural-language macro-goals, such
as ``buy cheaply, but do not lose the seller if the item seems scarce.'' Such
instructions combine soft preferences, situational conditions, and hard
constraints. Their meaning is also local. Early in a negotiation, the agent may
reasonably emphasize price gain; after the seller becomes firm, issues a final
offer, or signals walkaway risk, the same macro-goal may require a temporary
shift toward deal probability, rapport preservation, or turn efficiency. A
single static vector can describe one trade-off, but it does not explain how
the agent should change its short-term behavior as the conversation evolves.

Prior work on target-guided and proactive dialogue has made substantial
progress on multi-turn planning. Target-guided models learn keyword or topic
transitions toward a predefined target
\cite{tang-etal-2019-target,qin2020dynamic,zhong2021keyword,kishinami-etal-2022-target,yang-etal-2022-topkg},
while recent planners use Brownian-bridge trajectories, LLM reasoning,
self-reflection, experience retrieval, and dual-process planning to improve
future-oriented behavior
\cite{wang-etal-2023-dialogue,zheng-etal-2024-thoughts,guo-etal-2024-pcqpr,dao-etal-2024-experience,he2024planning}.
This literature shows that dialogue systems benefit from planning beyond the
next utterance. At the same time, much of it assumes that the target is fixed
before the interaction is optimized. The setting we study is different: the
macro-goal remains stable at the linguistic level, but the short-term skill
needed to pursue it can change as the other party's intent and stance change.

The closest multi-objective predecessor to our work is PADPP, which formulates
goal-oriented dialogue as preference-conditioned policy learning and trains a
single planner that adapts to different objective vectors at inference time
\cite{dao-liao-2025-one}. Its use of generalized policy improvement (GPI)
connects naturally to the broader line of work on reusable behavior across
reward functions \cite{barreto2017successor,barreto2018transfer}. PADPP
substantially reduces the need to retrain a separate policy for every objective
setting. However, it still assumes that the preference vector is given. This
leaves open a practical question at the interface between human intent and
policy control: how should an ambiguous macro-goal be decomposed into
short-term dialogue skills, when should the active skill change, and how should
that change be justified by evidence in the conversation?

This question becomes especially important under conversational intent drift.
Fine-grained intention understanding suggests that dialogue intent involves
multiple aspects, including situation, emotion, action, and knowledge
\cite{liang-etal-2025-intentionframe}, and recent drift detection work studies
how active intent changes over continuous conversations
\cite{wang2025detecting}. In our buyer-side negotiation setting, drift appears
as a change in the seller's stance: a seller may begin neutral, become firm
after repeated low offers, issue a final offer, or signal walkaway risk. If the
buyer continues executing the same short-term strategy, it may bargain past the
point at which the seller's stance has changed; if it switches too early, it
may sacrifice price utility unnecessarily. We therefore study goal-driven
dialogue policy learning in which a stable natural-language macro-goal must be
served through adaptive short-term skills under intent drift and hard
constraints.

We propose \textbf{ALMOND}, a hierarchical framework for learning adaptive
dialogue policies from natural-language macro-goals. ALMOND separates the
problem into two coupled levels. At the low level, the policy learns reusable
communication skills, such as bargaining, conceding, asking for clarification,
or preserving rapport. Because complex dialogue behaviors often depend on
skills that are unevenly learned, ALMOND uses a regret-based mechanism to give
more training attention to under-mastered skills rather than assuming that all
skills progress uniformly. At the high level, a policy monitors the dialogue
state for intent shifts and selects short-term skills conditioned on the
macro-goal, the recent interaction history, and strategic hints distilled
during training. In this design, language-level goal interpretation guides
strategic skill selection, while the learned low-level policy remains
responsible for executing concrete communicative behavior.

We evaluate ALMOND with metrics that reflect constrained adaptation rather
than final task completion alone. LLM-SR measures linguistic goal completion
using an LLM-as-Judge protocol with human-verification samples. Goal Success
Rate (GSR) counts a dialogue as successful only when the buyer reaches a deal
within the buyer's budget ceiling and turn limit. Turn-to-Drift Adaptation
Delay (T2DA) measures $t_{\mathrm{adapt}}-t_{\mathrm{drift}}$, and Constraint
Violation Rate (CVR) tracks blocked and executed safety violations. Across our
benchmark settings, ALMOND improves LLM-as-Judge goal completion by 4--7
percentage points and reduces adaptation delay by 2--3 turns over
static-weight and ablation baselines. Phase-1 skill training further improves
$r$-fair, $r$-gain, and $r$-deal by 3--6 percentage points while reducing
average turns.

Our contributions are:
\begin{itemize}
    \item We formulate goal-driven dialogue from ambiguous natural-language
    macro-goals under conversational intent drift and hard constraints.
    \item We propose ALMOND, a hierarchical framework that combines
    regret-aware low-level skill learning with drift-aware high-level skill
    selection.
    \item We use strategic hints distilled during training to guide short-term
    skill selection toward the overarching linguistic goal.
    \item We construct persona-conditioned seller simulators with deterministic
    drift triggers and evaluate constrained adaptation with LLM-SR, GSR, T2DA,
    and CVR.
\end{itemize}
